\chapter*{附录 D\quad 定理与定义索引}
\addcontentsline{toc}{chapter}{附录 D\quad 定理与定义索引}
\markboth{附录 D\quad 定理与定义索引}{}

\noindent 本索引按章编排，收录正文中全部带编号的定理、命题、引理、推论、定义、例与注等。
条目左列为正文编号，中列为条目主题，右列为该条目在原书日文版中的页码。
原书卷末索引以日文五十音排序，于中文检索不便，故不逐条移译，而以本索引代之。

\begingroup
\small
\begin{longtable}{@{}p{0.10\textwidth}p{0.10\textwidth}p{0.52\textwidth}p{0.10\textwidth}@{}}
\textbf{编号} & \textbf{类型} & \textbf{条目} & \textbf{原书页} \\
\hline
\endfirsthead
\textbf{编号} & \textbf{类型} & \textbf{条目} & \textbf{原书页} \\
\hline
\endhead
\multicolumn{4}{l}{\textbf{绪论}}\\*
0.1 & 定理 & $f\colon X\to Y$ 是概形的射，$F$ 是 $X$ 上的拟凝聚层。…… & 3 \\
0.2 & 定理 & [广中定理] $X$ 是 $\C$ 上的代数簇（或是解析空间的紧子集的足够小的邻域），$I\subset\cO_X$ 是定义闭子…… & 4 \\
0.3 & 定理 & [广中定理] $f\colon X\to C$ 是从 $\C$ 上的代数簇（或是解析空间的紧子集的足够小的邻域）到非奇异…… & 4 \\
0.4 & 记号 & 为避免误解，下面列举本书使用的标准记号。 & 4 \\
\multicolumn{4}{l}{\textbf{第 1 章\quad 有理曲线与标准类}}\\*
1.1 & 编号段 & 簇上的有理曲线，会对该簇的双有理几何产生怎样的影响？ & 10 \\
1.2 & 定理 & 设 $X$ 为固有的非奇异曲面。…… & 10 \\
1.3 & 命题 & 设 $X$ 为任意维数的非奇异簇，$f:Y\to X$ 为固有双有理态射，$x$ 为 $X$ 的任意点。…… & 10 \\
1.4 & 推论 & 设 $Z$ 为非奇异簇，$g: Z\dashrightarrow X$ 为有理映射。…… & 11 \\
1.5 & 推论 & 设 $X$、$Z$ 为代数簇，$Z$ 非奇异，$X$ 固有。…… & 11 \\
1.6 & 引理 & 刚性引理（rigidity lemma） & 11 \\
1.7 & 推论 & 折弯引理（bend and break）I & 12 \\
1.8 & 例 & 设 $E$ 为椭圆曲线，$M$ 为次数 $\ge 2$ 的线丛，它由截面 $s$、$t$ 生成。…… & 13 \\
1.9 & 引理 & 折弯引理 II & 14 \\
1.10 & 定理 & 设 $X$ 为非奇异射影簇，$-K_X$ 丰富。…… & 15 \\
1.11 & 原理 & 设整系数的齐次多项式系对无限多个素数 $p$ 在 $\bar{\mathbb{F}}_p$ 上有非平凡解（系数非整数时，改…… & 18 \\
1.12 & 编号段 & 下面减弱关于 $X$ 的假设。…… & 19 \\
1.13 & 定理 & 设 $X$ 为非奇异射影簇，$H$ 为 $X$ 上的丰富除子。…… & 20 \\
1.14 & 引理 & 设 $c,d>0$。…… & 21 \\
1.15 & 定义 & 设 $K=\Q$ 或 $K=\R$，$V$ 为 $K$-向量空间。…… & 22 \\
1.16 & 定义 & 设 $X$ 为固有簇。…… & 23 \\
1.17 & 定义 & 设 $X$ 为固有簇。…… & 24 \\
1.18 & 定理 & Kleiman 丰富性判别法（ampleness criterion） & 24 \\
1.19 & 推论 & 设 $X$ 为射影簇，$H$ 为丰富除子。…… & 24 \\
1.20 & 引理 & 设 $D$ 为固有不可约曲面上的除子，且 $(D^2)>0$。…… & 25 \\
1.21 & 推论 & 设 $X$ 为不可约射影曲面，$H$ 为丰富除子。…… & 26 \\
1.22 & 引理 & 设 $X$ 为不可约射影曲面，$C\subset X$ 为不可约曲线。…… & 26 \\
1.23 & 例 & 下面看几个例子。…… & 27 \\
1.24 & 定理 & 锥定理（cone theorem），[Mor82] & 28 \\
1.25 & 定义 & 设 $X$ 为射影簇，$F\subset\overline{NE}(X)$ 为极锥面。…… & 31 \\
1.26 & 注 & 如例 1.27 所见，一般来说并非任何极锥面都能收缩。…… & 32 \\
1.27 & 例 & 在非奇异平面三次曲线 $D$ 上取 $12$ 个点 $p_1,\dots,p_{12}$，把以这些点为中心的爆破记为 $g:X\to\PP^2$，令 $C$ 为 $D$ 的双有…… & 32 \\
1.28 & 定理 & 设 $X$ 为非奇异射影曲面，$R\subset\overline{NE}(X)$ 是满足 $(R\cdot K_X)<0$ 的极射线。…… & 33 \\
1.29 & 定理 & 设 $X$ 为非奇异射影曲面。…… & 34 \\
1.30 & 定义 & 若 $K_{X'}$ 是 nef（定理 1.29(1) 的情形），则称 $X'$ 为曲面 $X$ 的极小模型（minimal model）。…… & 34 \\
1.31 & 概述 & 至此所述，曲面极小模型的构成法给出如下步骤。 & 34 \\
1.32 & 定理 & 设 $X$ 为 $\C$ 上的三维非奇异射影簇，$\operatorname{cont}_R:X\to Y$ 为 $K_X$-负极射线 $R\subset\overline{NE}(X)$ 的收缩。…… & 35 \\
1.33 & 记号 & $X$ 是域 $k$ 上的固有概形，$Z\subset X$ 是 $d$ 维闭子概形，$L_1,\dots,L_d$ 是 $X$ 上的 Cartier 除子。…… & 37 \\
1.34 & 编号段 & （交点数的定义）\quad 定义交点数至少有下列 $4$ 种手法。 & 37 \\
1.35 & 命题 & 交点数具有下列性质。…… & 38 \\
1.36 & 定理 & 设 $X$ 是域上的 $n$ 维固有概形，$L_1,\dots,L_r$ 是 $X$ 上的 Cartier 除子。…… & 38 \\
1.37 & 定理 & 中井--Moishezon 判别法 & 39 \\
1.38 & 定理 & 设 $X$ 为固有簇，$L$ 为 nef Cartier 除子。…… & 40 \\
1.39 & 编号段 & （Kleiman 丰富性判别法（定理 1.18）的证明）\quad 设 $D$ 丰富。…… & 41 \\
1.40 & 注 & 当 $X$ 固有且非奇异（或至少是 $\Q$-分解的）时定理 1.18 也成立，但对一切固有簇并不成立。…… & 42 \\
1.41 & 命题 & 设 $f:X\to Y$ 为固有态射，$D$ 是 $X$ 上的 Cartier 除子，取一点 $y\in Y$，$X_y$ 是 $f$ 在 $y$ 上的纤维。…… & 42 \\
1.42 & 定理 & 相对的中井--Moishezon 判别法 & 44 \\
1.43 & 推论 & 设 $f:X\to Y$ 为固有态射，$L$ 是 $X$ 上的 f-nef Cartier 除子。…… & 44 \\
1.44 & 定理 & 相对 Kleiman 丰富性判别法 & 44 \\
1.45 & 命题 & 设 $f:X\to Y$ 为射影簇之间的态射，$M$ 是 $Y$ 上的丰富除子。…… & 44 \\
1.46 & 例 & 设 $0\in E$ 为椭圆曲线。…… & 44 \\
\multicolumn{4}{l}{\textbf{第 2 章\quad 极小模型纲领引论}}\\*
2.1 & 例 & 设 $A$ 为三维 Abel 簇，$X$ 为按二阶自同构 $x\mapsto -x$ 取 $A$ 的商所得的空间 $A/\pm$。…… & 48 \\
2.2 & 编号段 & 需要一种引入新概念的研究方法。…… & 48 \\
2.3 & 编号段 & 在概要 1.31(1) 中，首先考察 $K_X$ 是否为 nef。…… & 49 \\
2.4 & 编号段 & 在概要 1.31 步骤 2 中，需要收缩极射线 $R$。…… & 49 \\
2.5 & 命题 & 设 $X$ 为正规射影簇，并假设 $X$ 是 $\Q$-分解的（$\Q$-factorial）。…… & 49 \\
2.6 & 编号段 & 命题 2.5 的三种情形，从我们的纲领的立场看，行为大为不同。 & 50 \\
2.7 & 例 & （翻转的例子）用例子考察这一状况。…… & 51 \\
2.8 & 定义 & 设 $f\colon X\to Y$ 为固有双有理态射，且例外集 $\Exc(f)$ 在 $X$ 内的余维数至少为 2。…… & 53 \\
2.9 & 注 & 文献中对术语的用法并不统一。…… & 54 \\
2.10 & 编号段 & 在整理上述讨论的形式之前，必须先确定森纲领实际奏效的奇点类。…… & 54 \\
2.11 & 定义 & 设 $Y$ 为正规簇。…… & 54 \\
2.12 & 定义 & 设 $Y$ 为正规簇。…… & 55 \\
2.13 & 定义 & 设 $X$ 为正规且固有的簇。…… & 55 \\
2.14 & 编号段 & （森纲领，或极小模型纲领）\quad 从非奇异射影簇 $X$ 出发，反复进行易于理解的双有理变…… & 55 \\
2.15 & 例 & 这里给出若干高维极射线收缩的例子。…… & 57 \\
2.16 & 例 & 相对化 & 59 \\
2.17 & 例 & 解析的情形 & 60 \\
2.18 & 例 & 带群作用的簇 & 61 \\
2.19 & 定理 & 任意固有的三维 $G$-簇 $X$，都与某个 $GQ$-分解的、且只有终端奇点的 $G$-簇 $Y$ 连同 $G$-作用双有理等…… & 62 \\
2.20 & 例 & 对数簇 & 62 \\
2.21 & 例 & 相对化、解析的情形、带群作用的簇、对数簇 & 64 \\
2.22 & 定义 & 设 $X$ 为正规簇，且存在整数 $m>0$ 使 $mK_X$ 为 Cartier。…… & 64 \\
2.23 & 注 & 设 $k(X)$ 为 $X$ 的有理函数体。…… & 64 \\
2.24 & 定义 & 设 $X$ 为簇。…… & 65 \\
2.25 & 定义 & $(X,\Delta)$ 是正规簇 $X$ 与 $\Q$-除子 $\Delta$ 的对。…… & 65 \\
2.26 & 记号 & 用数值等价，除以 $m$ 后可写成 \[ K\_Y+f\_*^{-1}\Delta\equiv f^*(K\_X+\Delta)+\sum a(E\_i,X,\Delta)E\_i, \] 其中和式取遍一切例外的 $E_i$…… & 66 \\
2.27 & 引理 & 记号同上。…… & 66 \\
2.28 & 定义 & $(X,\Delta)$ 的差异系数（discrepancy）由下式给出…… & 66 \\
2.29 & 引理 & 设 $X$ 为非奇异簇，$\Delta=\sum a_iD_i$ 为互不相同素除子之和。…… & 67 \\
2.30 & 引理 & 设 $f\colon Y\to X$ 为正规簇之间的固有双有理态射。…… & 67 \\
2.31 & 推论 & 设 $\Delta=\sum a_iD_i$ 为 $\Q$-除子。…… & 67 \\
2.32 & 推论 & 设 $f\colon Y\to X$ 为奇点消解，$E\subset Y$ 为例外集，$E_i\subset E$ 取遍一切例外素除子。…… & 70 \\
2.33 & 推论 & 取对 $(X,\Delta)$ 与 $X$ 上的线性系 $|L|$。…… & 71 \\
2.34 & 定义 & 设 $(X,\Delta)$ 为正规簇 $X$ 与互不相同素除子之和 $\Delta=\sum a_iD_i$ 的对。…… & 72 \\
2.35 & 推论 & 设 $(X,\Delta)$ 为对，$\Delta'$ 为有效 $\Q$-Cartier 除子。…… & 73 \\
2.36 & 命题 & 设 $(X,\Delta)$ 为 klt 对。…… & 74 \\
2.37 & 定义 & 设 $(X,\Delta)$ 为正规簇 $X$ 与互不相同素除子之和 $\Delta=\sum a_iD_i$ 的对。…… & 75 \\
2.38 & 注 & 这个定义极为微妙。…… & 75 \\
2.39 & 推论 & 设 $(X,\Delta)$ 为 $\Delta$ 有效的对，$\Delta'$ 为有效 $\Q$-Cartier 除子。…… & 75 \\
2.40 & 命题 & 考虑正规簇 $X$ 与 $\Q$-除子 $\Delta$ 的对 $(X,\Delta)$。…… & 76 \\
2.41 & 命题 & 设 $(X,\Delta)$ 为 dlt 对。…… & 76 \\
2.42 & 命题 & 设 $(X,\Delta)$ 为曲面上的 dlt 对，$x\in X$。…… & 76 \\
2.43 & 命题 & 设 $(X,\Delta)$ 为 dlt 对，$X$ 为具有丰富除子 $H$ 的射影簇。…… & 77 \\
2.44 & 定理 & 设 $(X,\Delta=\sum d_iD_i)$，其中 $D_i$ 为互不相同的素除子，系数 $d_i$ 满足 $0\le d_i\le 1$。…… & 77 \\
2.45 & 引理 & 设 $X$ 为域 $k$ 上的代数簇。…… & 78 \\
2.46 & 原理 & 若层 $F$ 的上同调来自拓扑上同调，则存在小平型消没定理。 & 79 \\
2.47 & 定理 & 小平消没定理 & 80 \\
2.48 & 定理 & 设 $X,Y$ 为 $\C$ 上的有限型分离概形，$f:X\to Y$ 为固有态射，$F$ 为 $X$ 上的凝聚层。…… & 80 \\
2.49 & 定义 & 不分歧循环覆盖 & 80 \\
2.50 & 定义 & 分歧循环覆盖，线丛的情形 & 81 \\
2.51 & 引理 & $Z$ 为非奇异的充要条件是 $X$ 与 $D$ 都非奇异。 & 82 \\
2.52 & 定义 & 分歧循环覆盖，一般情形 & 82 \\
2.53 & 引理 & 记号同上。…… & 82 \\
2.54 & 推论 & 设 $U\subset X$ 为满足 $U\cap D\ne\varnothing$ 的连通开集。…… & 83 \\
2.55 & 引理 & 设 $F$ 为拓扑空间 $X$ 上 Abel 群的层，$F_1,F_2\subset F$ 为子层。…… & 83 \\
2.56 & 推论 & 设 $0<r<m$，$M\subset L^{-r}$ 为满足 $M|_{X\setminus D}=L^{-r}|_{X\setminus D}$ 的子层。…… & 84 \\
2.57 & 命题 & 设 $X$ 为 $\C$ 上的非奇异射影簇，$D,D'$ 为 $X$ 上的 Cartier 除子，且 $D\equiv D'$。…… & 84 \\
2.58 & 引理 & 设 $X$ 为域上 $n$ 维射影概形，$B$ 为 $X$ 上的 Cartier 除子。…… & 85 \\
2.59 & 定义 & 设 $X$ 为 $n$ 维固有簇。…… & 85 \\
2.60 & 引理 & 设 $X$ 为 $n$ 维射影簇，$D$ 为 Cartier 除子。…… & 85 \\
2.61 & 命题 & 设 $X$ 为 $n$ 维射影簇，$D$ 为 Cartier 除子。…… & 86 \\
2.62 & 引理 & 设 $f:Y\to X$ 为双有理态射。…… & 88 \\
2.63 & 推论 & 记号与假设同引理 2.62。…… & 88 \\
2.64 & 定理 & 小平消没定理 II，{[Kaw82, Vie82]} & 89 \\
2.65 & 主张 & 设 $X$ 为非奇异复簇，$L$ 为 $X$ 上的线丛，且满足 $L\equiv M+\sum a_iD_i$。…… & 89 \\
2.66 & 注 & 由 Serre 对偶，$H^i(X,L^{-1})$ 的消没定理与 $H^{n-i}(X,\omega_X\otimes L)$ 的消没定理等价…… & 91 \\
2.67 & 命题 & 设 $X$ 为拟射影簇，$D$ 为 $X$ 上的 Cartier 除子，$m$ 为自然数。…… & 91 \\
2.68 & 推论 & 设 $f:Y\to X$ 为射影簇之间的双有理态射，$Y$ 非奇异。…… & 92 \\
2.69 & 命题 & 设 $f:Y\to X$ 为固有态射，$F$ 为 $Y$ 上的凝聚层。…… & 92 \\
2.70 & 定理 & 一般小平消没定理 & 93 \\
2.71 & 命题 & 此时去掉 (2) 的 snc 条件，结论仍成立。 & 93 \\
\multicolumn{4}{l}{\textbf{第 3 章\quad 锥定理}}\\*
3.1 & 定理 & 设 $Y$ 是非奇异复射影簇，$\sum d_iD_i$ 是（写成不同素除子之和的）$\mathbb{Q}$-除子，$L$ 是线丛（或 Cartier 除子）。…… & 96 \\
3.2 & 编号段 & 陈述 4 个基本定理，最后一个便是锥定理。…… & 96 \\
3.3 & 定理 & 无基点定理 & 96 \\
3.4 & 定理 & 非消没定理 & 97 \\
3.5 & 定理 & 有理性定理 & 97 \\
3.6 & 补充 & （[KMM94a]）\quad 记号同上。…… & 97 \\
3.7 & 定理 & 锥定理 & 97 \\
3.8 & 编号段 & 这些定理的证明的逻辑顺序如下。 & 98 \\
3.9 & 编号段 & （基本方针）\quad 与非消没定理、有理性定理同样，无基点定理的证明的主要思想如下…… & 98 \\
3.10 & 编号段 & 进行以上讨论时，必须考虑把 $\mathbb{Q}$-除子及其取上整限制到非奇异簇 $Y$ 的非奇异素除…… & 100 \\
3.11 & 定理 & 设 $(X,\Delta)$ 是 $\Delta$ 有效的固有 klt 对。…… & 102 \\
3.12 & 猜想 & 丰度猜想 & 103 \\
3.13 & 编号段 & 如定理 3.11 中所示，若 $K_X+\Delta$ 是 nef 的，则 (2) 蕴含 (1)。…… & 103 \\
3.14 & 编号段 & （思路说明）\quad 说明有理性定理如何被用来获得关于曲线锥的信息。…… & 104 \\
3.15 & 定理 & 设 $N_{\mathbb{Z}}$ 是有限秩自由 $\mathbb{Z}$-模，$N_{\mathbb{R}}=N_{\mathbb{Z}}\otimes_{\mathbb{Z}}\mathbb{R}$ 是到 $\mathbb{R}$ 的系数扩张。…… & 104 \\
3.16 & 注 & 行文之便，这里暂且离开主线，先证明补充 3.6。…… & 106 \\
3.17 & 推论 & $(X,\Delta)$ 是射影 klt 对，$R\subset\overline{NE}(X)$ 是 $(K_X+\Delta)$-负的极射线，$g_R:X\to Z$ 是其收缩态射。…… & 108 \\
3.18 & 推论 & $(X,\Delta)$ 是射影 klt 对，$R\subset\overline{NE}(X)$ 是 $(K_X+\Delta)$-负的极射线，$g_R:X\to Z$ 是其收缩态射。…… & 108 \\
3.19 & 引理 & 给定正整数 $a$、正数 $\epsilon$ 与实数 $r$。…… & 110 \\
3.20 & 主张 & 无论怎样选取满足 $0<aq-rp<\epsilon$ 的充分大的 $(p,q)$，$L(p,q)$ 都是 $X$ 的同一子集。…… & 111 \\
3.21 & 编号段 & 下面是相对设定最原始的原型。 & 118 \\
3.22 & 定义 & $f:X\to f(X)\subset Y$ 是满足条件 3.21(0)--(2) 之一的固有态射，$X$ 不可约，$n$ 是 $f$ 的一般纤维的维数。…… & 118 \\
3.23 & 引理 & $f:X\to Y$ 是射影的，$X$ 不可约，$D$ 是 $X$ 上的 Cartier 除子。…… & 119 \\
3.24 & 定理 & 相对无基点定理（relative basepoint-free theorem） & 119 \\
3.25 & 定理 & 相对锥定理（relative cone theorem） & 119 \\
3.26 & 编号段 & （情形 3.21(0) 中定理 3.24 的证明）\quad 把 $D$ 换成 $D+f^*(\text{丰富除子})$，并把 $\Delta$ 稍作变化，在保持 $(X,\Delta)$ 是 klt 这一条件的同时，使 $D$ 成为 nef 的…… & 120 \\
3.27 & 命题 & $f:X\to Y$ 是射影簇的态射。…… & 121 \\
3.28 & 编号段 & （情形 3.21(0) 中定理 3.25 的证明）\quad 先证定理 3.25(1)。…… & 121 \\
3.29 & 编号段 & （对之类别的选取）\quad $X$ 是域上的正规概形或正规复解析流形。…… & 122 \\
3.30 & 编号段 & （对的双有理几何）\quad 一旦确定了应当使用的对之类，下一个必须直面的问题就是…… & 122 \\
3.31 & 编号段 & （$\Q$-分解的 dlt 对情形的相对 MMP）\quad 从 $X$ 为域上的正规概形或正规复解析流形的对 $(X,\Delta)=(X_0,\Delta_0)$ 出发。…… & 123 \\
3.32 & 编号段 & 本节余下部分说明纲领的哪些部分可由已确立的结果推出，哪些部分仍待证…… & 125 \\
3.33 & 定义 & 设 $X$ 是正规概形（或正规解析流形），$D$ 是 $X$ 上的 Q-除子，且 $K_X+D$ 为 Q-Cartier。…… & 125 \\
3.34 & 定义 & 固有双有理态射 $f:X\to Y$ 在满足下述条件时称为极端的（extremal）或射线型的（ray type）。…… & 126 \\
3.35 & 定理 & 相对锥定理 3.25 对 Q-分解的 dlt 对 $(X,\Delta)$ 与射影态射 $f:X\to S$ 成立。 & 127 \\
3.36 & 命题 & 设 $(X,\Delta)$ 为射影的 Q-分解 dlt 对，$g_R:X\to Y$ 为 $(K_X+\Delta)$-负的极射线的收缩。…… & 127 \\
3.37 & 命题 & 设 $(X,\Delta)$ 为射影的 Q-分解 dlt 对，$g_R:X\to Y$ 为 $(K_X+\Delta)$-负的极射线的翻转收缩，$g_R^{+}:X^{+}\to Y$ 为其翻转。…… & 127 \\
3.38 & 引理 & 考虑下述交换图。…… & 128 \\
3.39 & 引理 & （nef 例外除子引理） 设 $h:Z\to Y$ 是正规簇之间的固有双有理态射。…… & 129 \\
3.40 & 引理 & 设 $f:Y\to X$ 是正规曲面的奇点消解，$E_i$ 为例外曲线。…… & 130 \\
3.41 & 引理 & 设 $U$ 为非奇异曲面，$C=\cup C_i$ 是 $U$ 上固有曲线的集合。…… & 130 \\
3.42 & 推论 & 设 $(X,\Delta)$ 为对，$\varphi:X\dashrightarrow X^{+}$ 为 $(K_X+\Delta)$-翻转。…… & 131 \\
3.43 & 推论 & 设 $(X,\Delta)$ 为对，$f:X\to Y$ 为 $(K_X+\Delta)$-负的极射线的除子收缩，$E$ 为例外除子。…… & 131 \\
3.44 & 推论 & 设 $(X,\Delta)$ 为 dlt 对。…… & 132 \\
3.45 & 编号段 & （MMP 的终止）\quad 设 $(X_0,\Delta_0)$ 为 Q-分解 dlt 对，并假设运行 MMP 得到对 $(X_i,\Delta_i)$。…… & 132 \\
3.46 & 编号段 & 尚未解决的最重要问题是 $(K+\Delta)$-翻转的存在性。…… & 132 \\
3.47 & 定理 & 设 $(X,\Delta)$ 为二维 Q-分解射影 dlt 对。…… & 133 \\
3.48 & 定义 & 设 $(X,\Delta)$ 为对，$f:X\to S$ 为固有态射。…… & 133 \\
3.49 & 例 & 设 $S$ 是只有一个奇点 $P\in S$ 的射影曲面，$f:S'\to S$ 为其极小奇点消解。…… & 134 \\
3.50 & 定义 & 设 $X$ 为正规簇，$\Delta$ 为使 $K_X+\Delta$ 为 Q-Cartier 的（不必有效的）Q-除子，$(X,\Delta)$ 为对，$f:X\to S$ 为固有态射。…… & 134 \\
3.51 & 命题 & 在定义 3.50 中，假设 $(X^{w},\Delta^{w})$ 满足 (1)、(4)、(5) 与下列条件 [（$3'$）] $\varphi_*\Delta-\Delta^{w}$ 去掉 $\varphi^{-1}$-例外除子的部分后有效， 则下列条件满足…… & 136 \\
3.52 & 定理 & $(X,\Delta)$ 是 $K_X+\Delta$ 为 Q-Cartier 的对，$f:X\to S$ 为固有态射。…… & 136 \\
3.53 & 推论 & 设 $(X,\Delta)$ 与 $(X',\Delta')$ 是使 $K_X+\Delta$ 与 $K_{X'}+\Delta'$ 为 Q-Cartier 的对，$p:X'\to X$ 与 $f:X\to S$ 为固有态射。…… & 137 \\
3.54 & 推论 & 设 $f_i:X_i\to S$（$i=1,2$）为固有态射，$X_i$ 为终端，且 $K_{X_i}$ 为 $f_i$-nef。…… & 138 \\
\multicolumn{4}{l}{\textbf{第 4 章\quad 极小模型纲领的二维奇点}}\\*
4.1 & 记号 & 设 $X$ 是正规代数曲面或正规复解析曲面，$\Delta$ 是 $X$ 上的 $\Q$-除子。…… & 140 \\
4.2 & 推论 & 设 $U$ 是非奇异曲面，$C=\bigcup C_i$ 是 $U$ 上的连通固有曲线，且相交矩阵 $(C_i\cdot C_j)$ 负定。…… & 141 \\
4.3 & 推论 & 若 $K_Y+f_*^{-1}\Delta$ 是 $f$-nef 的，则或者对一切 $i$ 有 $a_i=0$，或者对一切 $i$ 有 $a_i<0$。 & 141 \\
4.4 & 定义 & 设 $(0\in X)$ 是正规曲面奇点的芽，$f\colon Y\to X$ 是极小奇点消解。…… & 141 \\
4.5 & 定理 & 记号同 4.1。…… & 141 \\
4.6 & 定义 & 设 $C=\bigcup C_i$ 是非奇异曲面 $U$ 上的一族固有曲线。…… & 142 \\
4.7 & 定理 & 设 $(0\in X,\Delta)$ 是数值 lc 对，$f\colon Y\to X$ 是极小奇点消解，$\Gamma$ 是 $f$-例外曲线 $E=\bigcup E_i$ 的对偶图。…… & 143 \\
4.8 & 注 & （关于术语） (1) 型奇点中 $E=E_1$ 非奇异者称为单纯椭圆型（simple elliptic）。…… & 145 \\
4.9 & 注 & 定理 4.7 的逆也成立。…… & 146 \\
4.10 & 命题 & 设 $Y$ 是非奇异射影曲面，$E_i\subset Y$ 是不可约曲线，全体构成 snc 除子 $\bigcup E_i$。…… & 146 \\
4.11 & 命题 & 设 $(X,\Delta)$ 是数值 dlt 对。…… & 147 \\
4.12 & 定理 & 设 $(0\in X,\Delta)$ 是数值 dlt 曲面对，$f\colon Y\to X$ 是极小奇点消解。…… & 148 \\
4.13 & 引理 & 设 $U$ 是（不必约化的）解析空间，且 $H^1(U,\cO_U)=0$ 成立。…… & 148 \\
4.14 & 推论 & 设 $f\colon Y\to(P\in X)$ 是曲面奇点芽的奇点消解，且 $R^1f_*\cO_Y=0$…… & 149 \\
4.15 & 定理 & 有一个成立。…… & 150 \\
4.16 & 定理 & 设 $(0\in X,\Delta)$ 是对数标准对，$f\colon Y\to X$ 是极小奇点消解，$\Gamma$ 是例外曲线 $E=\bigcup E_i$ 的对偶图。…… & 151 \\
4.17 & 定理 & 设 $(0\in X)$ 是正规曲面奇点，$f\colon Y\to X$ 是极小奇点消解，$\Gamma$ 是 $f$-例外曲线 $E=\bigcup E_i$ 的对偶图。…… & 152 \\
4.18 & 命题 & 对正规解析曲面的芽 $(0\in X)$，下列条件等价。…… & 152 \\
4.19 & 编号段 & 本节中的证明，除注 4.9、命题 4.10、定理 4.12、命题 4.18 这 4 个例外，在任意特征下都成立。…… & 153 \\
4.20 & 定理 & 设 $(0\in X)$ 是正规曲面奇点的芽。…… & 153 \\
4.21 & 注 & Du Val 奇点在许多不同的情形中自然出现，并且有许多种刻画方法 [Dur79]。…… & 153 \\
4.22 & 定理 & Du Val 奇点的嵌入维数（embedding dimension），即把 $m$ 取为极大理想时 $m/m^2$ 作为向量空间的维数，总是 $3$。…… & 155 \\
4.23 & 引理 & 设 $(0\in X)$ 是 Du Val 奇点。…… & 156 \\
4.24 & 编号段 & 这里反复使用 $4$ 种方法。 & 157 \\
4.25 & 编号段 & 证明分若干阶段进行，依次把方程归结为更简单的形式。 & 158 \\
4.26 & 定义 & 设 $f\colon X\to S$ 是概形或解析空间的态射。…… & 160 \\
4.27 & 例 & 设 $f\colon X\to C$ 是到非奇异曲线的平坦态射，且存在点 $0\in C$ 使 $f$ 在 $C\setminus\{0\}$ 上光滑。…… & 161 \\
4.28 & 定理 & 设 $f\colon (x\in X)\to(0\in S)$ 是解析流形的芽之间的态射，$X_0$ 是在 $x$ 处有 Du Val 奇点的曲面。…… & 161 \\
4.29 & 定义 & 设 $0\in X_0=(f(x_1,\dots,x_m)=0)\subset\C^m$ 是原点处孤立超曲面奇点的芽。…… & 162 \\
4.30 & 例 & 通过具体计算可知 $\tau(A_n)=n$、$\tau(D_n)=n$、$\tau(E_n)=n$ 成立。 & 162 \\
4.31 & 命题 & 设 $0\in X_0\subset\C^3$ 是 Du Val 奇点。…… & 162 \\
4.32 & 构造 & 设 $Z$ 是非奇异曲面，$C\subset Z$ 是非奇异曲线。…… & 163 \\
4.33 & 引理 & 此时下列结论成立。…… & 164 \\
4.34 & 构造 & 设 $Z$ 是非奇异曲面，$C\subset Z$ 是非奇异曲线。…… & 164 \\
4.35 & 引理 & 此时 $(P_1,\dots,P_n)$ 上 $Y_n\to C^n$ 的纤维是构成 4.32 中定义的 $Z$ 在 $P_1,\dots,P_n\in C$ 处的爆破。 & 164 \\
4.36 & 例 & （$A_n$-型）设 $0\in\PP^2$ 是一点，$Z=B_0\PP^2$，$q\colon B_0\PP^2\to\PP^2$ 为射影。…… & 164 \\
4.37 & 例 & （$D_n$-型）设 $Z=\PP^1\times\PP^1$，$q\colon Z\to\PP^1$ 为第 2 射影。…… & 164 \\
4.38 & 例 & （$E_n$-型）设 $Z=\PP^2$，$C\subset Z$ 是非奇异 3 次曲线，$P_0\in C$ 是其拐点（flex）。…… & 165 \\
4.39 & 定理 & 上述构成具有下列性质。…… & 165 \\
4.40 & 注 & (1) $Z_n(0)$ 上曲线的形状由下列图表示。…… & 166 \\
4.41 & 编号段 & 定理 4.28 的证明 & 169 \\
4.42 & 注 & 由前述 Du Val 奇点形变的具体描述，可以简单地刻画邻近纤维中可能出现的奇点。…… & 169 \\
4.43 & 定理 & （[Bri71, Art74]）设 $f\colon X\to S$ 是曲面的平坦族，使每个纤维 $X_s$ 都只含 Du Val 奇点。…… & 169 \\
4.44 & 记号 & 本节中，$(0\in X)$ 表示 $0$ 是唯一奇点的正规曲面芽。…… & 170 \\
4.45 & 命题 & 记号同上。…… & 170 \\
4.46 & 引理 & 记号同上。…… & 171 \\
4.47 & 命题 & 记号同 4.44。…… & 171 \\
4.48 & 定义 & 当 $K_X$ 是 Cartier 的且 $R^1f_*\cO_Y\cong\C$ 成立时，称 $(0\in X)$ 为椭圆奇点（elliptic singularity）。…… & 172 \\
4.49 & 引理 & 假设 $(0\in X)$ 是椭圆的，$Z$ 如记号 4.44。…… & 172 \\
4.50 & 引理 & 设 $V$ 是（不必约化的）固有曲线且满足 $H^1(\cO_V)=0$，$L$ 是 $V$ 上的 nef 线丛。…… & 173 \\
4.51 & 命题 & $(0\in X)$ 是椭圆的，$Z$ 如记号 4.44。…… & 173 \\
4.52 & 记号 & $(0\in X)$ 是二维椭圆奇点，$L$ 是 $Z$ 上满足 $(L\cdot Z)>0$ 的 nef 线丛，截面 $s\in H^0(Z,L)$ 选取为使 $\{s=0\}$ 是与 $\operatorname{red}(Z)$ 的奇点集合不相交的零…… & 174 \\
4.53 & 引理 & 记号与假设同 4.52。…… & 174 \\
4.54 & 命题 & （[Lau77, Rei76]）$(0\in X)$ 是二维椭圆奇点，$Z$ 如记号 4.44。…… & 175 \\
4.55 & 引理 & 设 $R=\bigoplus_{i\ge 0}R(i)$ 是分次 $\C$-代数，$R(0)=\C$，$s\in R(1)$ 是非零因子。…… & 178 \\
4.56 & 定义 & 首先固定 $w_1,\dots,w_n\in\Z_{>0}$。…… & 178 \\
4.57 & 定理 & （[Lau77, Rei76]）设 $(0\in X)$ 是二维椭圆奇点，$f:Y\to X$ 是极小奇点消解。…… & 179 \\
4.58 & 注 & 上面由权定义的分次依赖于坐标的选取，但理想 $m^w(n)$ 与坐标的选取无关（参见断言…… & 179 \\
4.59 & 定义 & 设 $0\in X_0$ 是复解析空间的芽。…… & 181 \\
4.60 & 定理 & （[Gra72]）对任意孤立奇点，半普遍形变空间存在，且在局部解析同构的意义下唯一确定…… & 182 \\
4.61 & 定理 & 设 $0\in X_0=(f(x_1,\dots,x_m)=0)\subset\C^m$ 是原点处孤立超曲面奇点的芽。…… & 182 \\
4.62 & 定理 & 设 $(0\in T)$ 是复解析空间的芽，$F(x,t)\in\cO_{0,\C^m\times T}$ 是满足 $F(x,0)=f(x)$ 的收敛幂级数。…… & 183 \\
4.63 & 引理 & 设 $J\subset\cO_{0,\C^m\times T}$ 是使 $\cO_{0,\C^m\times T}/J$ 在 $\cO_{0,T}$ 上有限的理想。…… & 183 \\
4.64 & 引理 & 设 $I=(F,\partial F/\partial x_1,\dots,\partial F/\partial x_m)$，$I_1\subset(x,t)\cO_{0,\C^m\times T}$ 是任意理想，$G(x,t)\in\cO_{0,\C^m\times T}$ 是满足 \[ F(x,t)\equiv G(x,t)\mod I^2\cdot I\_1 \] 的任意元。…… & 184 \\
4.65 & 编号段 & 定理 4.62 的证明 & 185 \\
4.66 & 编号段 & 代数情形下的结果 & 186 \\
4.67 & 定义 & 设 $f:(\xi\in X)\to(\eta\in Y)$ 是 $\C$ 上代数概型芽之间的态射。…… & 187 \\
4.68 & 推论 & （引理 4.64 及其代数类似的推论）设给定两个在原点处有孤立奇点的多项式超曲面…… & 187 \\
4.69 & 命题 & 设 $f,g_1,\dots,g_n$ 是正则有理函数。…… & 187 \\
4.70 & 命题 & 设 $f:(x\in X)\to(t\in T)$ 是 $\C$ 上代数概型芽的态射，$d=\dim_xf^{-1}(t)$。…… & 189 \\
4.71 & 推论 & （极小嵌入（minimal embedding））符号同上。…… & 189 \\
4.72 & 引理 & 设 $(x\in X)\to(y\in Y)$ 是 $\C$ 上代数概型芽的拟有限态射。…… & 190 \\
\multicolumn{4}{l}{\textbf{第 5 章\quad 极小模型纲领的奇点}}\\*
5.1 & 定义 & 设 $(R,\mathfrak{m})$ 是 Noether 局部环，$N$ 是有限 $R$-模。…… & 192 \\
5.2 & 定义 & Noether 概形 $X$ 上的层 $F$，若对每个 $x\in X$，茎 $F_x$ 具有长度至少为 $\min\{d,\dim\cO_{X,x}\}$ 的正则列，则称 $F$ 为 $S_d$。…… & 192 \\
5.3 & 命题 & 设 $(R,\mathfrak{m})$ 是 Noether 局部环，$N$ 是有限 $R$-模。…… & 193 \\
5.4 & 命题 & 设 $f\colon X\to Y$ 是 Noether 概形之间的有限态射，$F$ 是 $X$ 上的凝聚层，且 $\Supp F$ 为纯 $r$ 维。…… & 193 \\
5.5 & 推论 & 设 $f\colon X\to Y$ 是纯 $n$ 维概形之间的有限态射，$F$ 是 $X$ 上的凝聚层。…… & 194 \\
5.6 & 定义 & 设 $X$ 是概形，$A$ 是 $\cO_X$-代数的局部自由层。…… & 194 \\
5.7 & 命题 & 设 $f\colon X\to Y$ 是纯 $n$ 维概形之间的有限满射。…… & 194 \\
5.8 & 定义 & 设 $Y$ 是特征 $0$ 的域上的簇，$f\colon X\to Y$ 是奇点消解。…… & 194 \\
5.9 & 编号段 & 有理奇点的概念在正特征下运作得并不好。…… & 195 \\
5.10 & 定理 & 设 $Y$ 是特征 $0$ 的域上的簇。…… & 195 \\
5.11 & 推论 & 若 $X$ 是 $\C$ 上的簇，则 \[ X\ \text{具有有理奇点}\iff X^{an}\ \text{具有有理奇点}. \] & 196 \\
5.12 & 引理 & 设 $Y$ 是特征 $0$ 的域 $k$ 上 $n$ 维簇，$f\colon X\to Y$ 是奇点消解。…… & 196 \\
5.13 & 命题 & 设 $f\colon X\to Y$ 是特征 $0$ 的域上 $n$ 维簇之间的有限态射。…… & 198 \\
5.14 & 定义 & 设 $x\in X$ 是复解析空间的芽。…… & 198 \\
5.15 & 命题 & 设 $X$ 是 $\C$ 上的代数簇或解析空间，且只有商奇点。…… & 198 \\
5.16 & 引理 & 设 $f\colon X\to Y$ 是正规簇之间的有限满射。…… & 199 \\
5.17 & 引理 & 设 $(X,\Delta)$ 是对，$|H|$ 是 $X$ 上的无基点线性系，$H_g\in|H|$ 是一般元。…… & 199 \\
5.18 & 推论 & 若 $X$ 终端，则 $X$ 在余维 $2$ 处非奇异。…… & 199 \\
5.19 & 定义 & 设 $X$ 是正规簇，$D$ 是 $X$ 上的 $\Q$-Cartier Weil 除子。…… & 200 \\
5.20 & 命题 & 设 $g\colon X'\to X$ 是 $n$ 维正规簇之间的有限态射。…… & 200 \\
5.21 & 推论 & 设 $x\in X$ 是正规奇点的芽。…… & 202 \\
5.22 & 定理 & 设 $(X,\Delta)$ 是 dlt 对，且 $\Delta$ 有效。…… & 202 \\
5.23 & 例 & 维数至少为 $2$ 的 Abel 簇上的锥，其顶点为 lc，但并非 CM。 & 204 \\
5.24 & 推论 & 设 $X$ 是使 $K_X$ 为 Cartier 的正规簇。…… & 204 \\
5.25 & 推论 & 设 $(X,\Delta)$ 是 dlt 对，$D$ 是 $X$ 上的 $\Q$-Cartier Weil 除子。…… & 204 \\
5.26 & 命题 & 设 $(X,\Delta)$ 是射影的 dlt 对，$S,D$ 是 $\Q$-Cartier Weil 除子。…… & 204 \\
5.27 & 推论 & （dlt 对的 Serre 对偶）设 $(X,\Delta)$ 是特征 $0$ 的域上的纯 $n$ 维射影 dlt 对，$D$ 是 $X$ 上的 $\Q$-Cartier Weil 除子。…… & 205 \\
5.28 & 定理 & 设 $f\colon X\to S$ 是平坦态射，纤维只有标准奇点。…… & 205 \\
5.29 & 记号 & 本节的讨论在代数、解析两种设定下都有效。…… & 206 \\
5.30 & 引理 & 设 $(0\in X)$ 是指数 $1$ 的三维标准奇点，$0\in H\subset X$ 是过 $0$ 的一般超曲面截面。…… & 206 \\
5.31 & 注 & 上述证明在所有维数下都有效。…… & 207 \\
5.32 & 定义 & 设 $(0\in X)$ 是三维奇点。…… & 207 \\
5.33 & 定义 & 设 $(0\in X)$ 是复解析奇点。…… & 207 \\
5.34 & 定理 & 设 $(0\in X)$ 是三维奇点，$0\in H\subset X$ 是过 $0$ 的一般超曲面截面。…… & 207 \\
5.35 & 定理 & 设 $0\in X$ 是指数 $1$ 的三维标准奇点。…… & 208 \\
5.36 & 编号段 & （定理 5.34 与定理 5.35 的证明方针）\quad 如下同时证明这两个定理。 & 208 \\
5.37 & 主张 & 设 $f(x_1,x_2,x_3,x_4)\in\C[[x_1,x_2,x_3,x_4]]$ 是幂级数。…… & 211 \\
5.38 & 推论 & 设 $0\in X$ 是三维正规奇点。…… & 212 \\
5.39 & 推论 & 若 $0\in X$ 是三维终端奇点，则它是孤立 cDV 奇点（或非奇异点）$0\in Y$ 由在原点之外自由作用的循…… & 212 \\
5.40 & 推论 & 设 $X$ 是三维簇，且只有至多标准奇点。…… & 212 \\
5.41 & 推论 & 设 $0\in X'$ 是三维标准奇点，其指数为 $r$。…… & 212 \\
5.42 & 定理 & 设 $0\in X$ 是奇点，$0\in H\subset X$ 是 Cartier 除子。…… & 213 \\
5.43 & 定理 & 设 $(0\in X)$ 是三维正规孤立奇点。…… & 214 \\
5.44 & 推论 & 设 $f\colon X\to S$ 是纤维维数为 $3$ 的平坦态射。…… & 215 \\
5.45 & 定义 & 设 $(X,\Delta)$ 是对，$Z\subset S\subset X$ 是闭部分概形。…… & 215 \\
5.46 & 命题 & 设 $X$ 是正规簇，$S$ 是余维 $2$ 且为 Cartier 的正规 Weil 除子，$Z\subset S$ 是闭子簇，$B=\sum b_iB_i$ 是 $\Q$-除子。…… & 216 \\
5.47 & 注 & 本节反复使用的假定是“$S\subset X$ 为余维 $2$ 且 Cartier 的 Weil 除子”。…… & 216 \\
5.48 & 定理 & 连通性定理 [K+92, 17.4] & 217 \\
5.49 & 推论 & 设 $g\colon Y\to X$ 是正规簇之间的固有双有理态射，$D=\sum d_iD_i$ 是 $Y$ 上的 $\Q$-除子，且 $g_*D$ 有效，$-(K_Y+D)$ 为 $g$-nef 的 $\Q$-Cartier。…… & 218 \\
5.50 & 定理 & 伴随反转 & 218 \\
5.51 & 命题 & 设 $(X,\Delta)$ 是 dlt 对，$\Delta$ 有效。…… & 218 \\
5.52 & 推论 & 设 $(X,\Delta)$ 是 dlt 对，$\Delta$ 有效。…… & 219 \\
5.53 & 推论 & 设 $(X,S+\Delta)$ 是 dlt 对，$\Delta$ 有效。…… & 219 \\
5.54 & 注 & 推论 5.53 对 lc 对 $(X,S+\Delta)$ 不成立。…… & 220 \\
5.55 & 推论 & 设 $(0\in X,\Delta)$ 是二维 dlt 对的芽，$\Delta$ 有效。…… & 220 \\
5.56 & 推论 & 设 $(S,D+\Delta)$ 是 dlt 对，$D$ 为 Cartier 除子，$\Delta$ 为有效 $\Q$-除子。…… & 220 \\
5.57 & 推论 & 设 $(0\in X,D+\Delta)$ 是 $0\in X$ 与 $0\in D$ 均非奇异、$\Delta$ 有效的二维对，且 $D$ 与 $\Delta$ 无公共成分。…… & 220 \\
5.58 & 命题 & 设 $f\colon S\to T$ 是正规曲面芽 $(0\in T)$ 的奇点消解。…… & 220 \\
5.59 & 命题 & 设 $X$ 是正规簇，$S$ 是余维 $2$ 且为 Cartier 的既约 Weil 除子，$B=\sum b_iB_i$ 是有效 $\Q$-除子。…… & 222 \\
5.60 & 注 & 命题 5.59 的逆不成立。…… & 222 \\
5.61 & 命题 & 设 $X$ 是正规簇，$S$ 是余维 $2$ 且为 Cartier 的既约 Weil 除子，$B=\sum b_iB_i$ 是有效 $\Q$-除子。…… & 222 \\
5.62 & 推论 & 设 $X$ 是正规簇，$S$ 是余维 $2$ 且为 Cartier 的既约 Weil 除子，$S'$ 为既约 Weil 除子，$B$ 为有效 $\Q$-除子。…… & 222 \\
5.63 & 命题 & 设 $X$ 是正规簇，$S_i$ 是过 $0\in X$ 的 $\Q$-Cartier Weil 除子。…… & 223 \\
5.64 & 注 & 命题 5.63 是下述结果 [K+92, 18.2] 的特殊情形：若 $(X,\sum b_iS_i)$ 为 lc，则 $\sum b_i\le\dim X$。 & 223 \\
5.65 & 编号段 & （$\PP^n$ 上的 Serre 对偶）\quad 设 $F$ 是（域 $k$ 上的）$\PP^n$ 上的凝聚层。…… & 224 \\
5.66 & 定义 & 设 $X$ 是域 $k$ 上的 $n$ 维固有概形。…… & 224 \\
5.67 & 命题 & 设 $f\colon X\to Y$ 是有限态射，$F,G$ 分别是 $X,Y$ 上的凝聚层。…… & 225 \\
5.68 & 命题 & 设 $f\colon X\to Y$ 是纯 $n$ 维固有概形之间的有限态射。…… & 225 \\
5.69 & 推论 & 对域 $k$ 上的任意纯 $n$ 维射影概形 $X$，$\omega_X$ 存在且为 $S_2$。 & 225 \\
5.70 & 推论 & 设 $X$ 是域 $k$ 上的纯 $n$ 维射影概形，$F$ 是 $X$ 上的凝聚层，$\Supp F$ 为纯 $n$ 维。…… & 225 \\
5.71 & 定理 & CM 层的 Serre 对偶（Serre duality） & 226 \\
5.72 & 推论 & 设 $X$ 是域 $k$ 上的纯 $n$ 维射影概形，$D$ 为丰富 Cartier 除子。…… & 226 \\
5.73 & 命题 & 伴随公式（adjunction formula） & 227 \\
5.74 & 注 & 即使 $X$ 不为 CM，由上述证明也可得到单射 $\omega_X(D)\otimes\cO_D\to\omega_D$。…… & 228 \\
5.75 & 命题 & 设 $X$ 是域 $k$ 上的 $n$ 维正规射影簇。…… & 229 \\
5.76 & 主张 & $\Phi$ 延拓为（用同样记号表示）下述双线性映射 & 229 \\
5.77 & 命题 & 相对迹映射（relative trace map） & 230 \\
5.78 & 引理 & 设 $K$ 为域，$f\in K[x]$ 为无重根的 $r$ 次首一（monic，即最高次系数为 $1$）多项式。…… & 230 \\
\multicolumn{4}{l}{\textbf{第 6 章\quad 三维翻模}}\\*
6.1 & 注 & 本节中，所有的定义与定理等，都是对正规簇叙述的。…… & 234 \\
6.2 & 引理 & 设 $Y$ 为正规代数簇（或解析空间），$B$ 为 $Y$ 上的 Weil 除子。…… & 234 \\
6.3 & 注 & 设 $Y=\Spec S$ 为仿射簇，$B$ 是由理想 $I\subset S$ 定义的有效除子。…… & 235 \\
6.4 & 推论 & 设 $f:X\to Y$ 为 $(K+D)$-翻转收缩。…… & 235 \\
6.5 & 定义 & 设 $f:X\to Y$ 为射线型翻转收缩（定义 3.34）。…… & 236 \\
6.6 & 命题 & 设 $Y$ 为正规簇，$B$ 为 $Y$ 上的 Weil 除子。…… & 236 \\
6.7 & 推论 & 设 $X,Y$ 为正规簇或 $\C$ 上的复解析空间，$f:X\to Y$ 为 $(K_X+D)$-翻转收缩。…… & 237 \\
6.8 & 命题 & （[Kaw88, 3.2]）设 $Y$ 与 $Y'$ 为正规且不可约的簇，$h:Y'\to Y$ 为有限满射。…… & 237 \\
6.9 & 推论 & 设 $f:X\to Y$ 为 $(K+D)$-翻转收缩。…… & 238 \\
6.10 & 定义 & 设 $X$ 为正规簇。…… & 238 \\
6.11 & 引理 & 设 $f:X\to Y$ 为翻转收缩。…… & 239 \\
6.12 & 命题 & 设 $f:X\to Y$ 为 $(K+D)$-翻转收缩。…… & 239 \\
6.13 & 注 & 本节的讨论，在代数的设定与解析的设定下同样有效。…… & 239 \\
6.14 & 定理 & 三维终端翻模的存在性 & 240 \\
6.15 & 定理 & 三维终端翻模保持解析奇点。…… & 241 \\
6.16 & 定义 & 所谓 $(K_X+\Delta)$-翻转的列（sequence of $(K_X+\Delta)$-flips），是指对 $(X^i,\Delta^i)$ 与双有理映射 $\varphi^i:X^i\dashrightarrow X^{i+1}$ 的列：它从对 $(X^1,\Delta^1)=(X,\Delta)$ 开始，其中 $\varphi^i:X^i\dashrightarrow X^{i+1}$ 是 $(K_{X^i}+\Delta^i)$-翻转，且 $\Delta^{i+1}:=\varphi^i_*\Delta^i$ 依次确定…… & 241 \\
6.17 & 定理 & 三维标准翻转的终止 & 241 \\
6.18 & 注 & 若 $X$ 是解析空间，则在定理 6.17 与推论 6.19 中须假设 $X$ 是紧致子集的充分小邻域。…… & 241 \\
6.19 & 推论 & 三维终端翻模的终止 & 241 \\
6.20 & 定义 & 设 $(X,\Delta=\sum a_iD_i)$ 是标准对，$D_i$ 是互不相同的素除子。…… & 241 \\
6.21 & 引理 & 设 $\varphi:X\dashrightarrow X'$ 是三维标准对 $(X,\Delta=\sum_{i=1}^k a_iD_i)$ 的 $(K_X+\Delta)$-翻转。…… & 242 \\
6.22 & 定义 & 设 $(X,\Delta)$ 是对，$f:Y\to X$ 是双有理态射。…… & 243 \\
6.23 & 定理 & 终端化（terminalization） & 243 \\
6.24 & 注 & 以 $\pi_X(:=f):X^{ter}(:=Y)\to X$ 表示定理 6.23 的证明中构造的终端化，则它具有如下性质。…… & 243 \\
6.25 & 定理 & Q-分解（Q-factorialization） & 243 \\
6.26 & 注 & 关于定理 6.25 所给出的 Q-分解 $f:Y\to X$，指出以下几点。…… & 243 \\
6.27 & 定理 & 设 $X$ 是只具有标准奇点的三维簇，且对每个奇点 $P\in X$，$(P\in X)$ 的指数 1 覆盖是 cDV 奇点。…… & 244 \\
6.28 & 引理 & 设 $f:U\to V$ 是固有双有理态射，$D$ 是 $U$ 上的 Cartier 除子。…… & 246 \\
6.29 & 注 & 若定理 6.27 中的 $X$ 是代数簇，则证明中只需用到射影簇的 MMP。…… & 246 \\
6.30 & 定义 & 设 $X$ 是至多具有标准奇点的三维簇。…… & 247 \\
6.31 & 引理 & 设 $\pi:(\tilde P\in\tilde X)\to(P\in X)$ 是标准奇点的任意次指数 1 覆盖（定义 5.19）。…… & 248 \\
6.32 & 定理 & 设 $X$ 是至多具有终端奇点的三维代数簇（或解析空间），$D$ 是 Weil 除子。…… & 248 \\
6.33 & 注 & 定理 6.32 所给出的 $g:Y\to X$（称为由 $X$ 的 $-D$ 决定的符号爆破（symbolic blowup））具有与注意 6.24 中所述基准终端化同样…… & 248 \\
6.34 & 注 & 若定理 6.32 中的 $X$ 是代数簇，则与注意 6.29 一样，可以只用射影簇的 MMP 给出证明。…… & 249 \\
6.35 & 引理 & 设 $h:Z\to X$ 是 $\C$ 上的至多具有有理奇点的簇的奇点消解，$x\in X$ 是使限制映射 $\operatorname{Pic}Z\to\operatorname{Pic}h^{-1}(x)$ 为满射的点。…… & 250 \\
6.36 & 引理 & 设 $X$ 只具有终端奇点，$Z\subset X$ 是余维 $\ge2$ 的闭子集。…… & 250 \\
6.37 & 注 & 本节的讨论（只要使用相应设定下的锥定理与 MMP 等）在射影的、非固有代数的以及…… & 251 \\
6.38 & 定理 & 设给定射影态射 $f_i:X_i\to Z$（$i=1,2$），以及与之相容的 $Z$ 上的双有理映射 $\pi:X_1\dashrightarrow X_2$。…… & 251 \\
6.39 & 引理 & 设 $\pi:X_1\dashrightarrow X_2$ 是 $Z$ 上的双有理映射，且 $\pi$ 在余维 1 处是同构。…… & 251 \\
6.40 & 命题 & 设 $f:X\to Z$ 是射影态射，$D$ 是 $X$ 上的有效除子且 $-D$ 是 $f$-丰富的。…… & 252 \\
6.41 & 推论 & 设 $Z$ 是至多具有标准奇点的射影簇或 Stein 簇，$f:X\to Z$ 是 crepant 射影双有理态射。…… & 252 \\
6.42 & 定理 & 设 $Z$ 是至多具有标准奇点的三维簇。…… & 253 \\
6.43 & 定理 & 三维标准翻模的终止定理（简版） & 253 \\
6.44 & 定理 & 三维标准翻模的存在 {[Kaw88, Kol89]} & 254 \\
6.45 & 推论 & 设 $X$ 是至多具有标准奇点的三维簇，$f:X\to Y$ 是（不必射线型的）$D$-翻模收缩。…… & 255 \\
6.46 & 推论 & 三维终端翻模的终止 & 256 \\
6.47 & 定理 & 三维标准翻模的终止 {[Kol89]} & 256 \\
\multicolumn{4}{l}{\textbf{第 7 章\quad 半稳定极小模型}}\\*
7.1 & 定义 & 设 $X$ 是正规簇，$B=\sum b_iB_i$ 是有效 $\Q$-除子，且 $f\colon X\to C$ 是到非奇异曲线 $C$ 的非常数态射。…… & 258 \\
7.2 & 引理 & 假设射 $f\colon(X,B)\to C$ 是 lc 的。…… & 258 \\
7.3 & 注 & 尽管引理 7.2(4) 给人以某种印象，$(X,B)$ 一般并不是标准的。…… & 259 \\
7.4 & 引理 & 假设 $f\colon(X,B)\to C$ 是 lc 射，$c\in C$，$x\in f^{-1}(c)$ 是闭点。…… & 259 \\
7.5 & 注 & 许多应用中，理解 $f^{-1}(c)$ 的非正规奇点也很重要。…… & 260 \\
7.6 & 引理 & 假设射 $f\colon(X,B)\to C$ 是 lc 的。…… & 260 \\
7.7 & 定理 & 设 $(X,B)$ 是 Q-分解的三维 dlt 对，$B$ 有效，$f\colon X\to Y$ 是射影态射。…… & 260 \\
7.8 & 定理 & （三维半稳定翻转的存在（参见推论 7.42））设 $(X,B)$ 是三维 Q-分解 dlt 对，$B$ 有效。…… & 261 \\
7.9 & 定理 & （半稳定极小模型定理）设给定三维 Q-分解对 $(X,B)$ 与非奇异曲线 $C$。…… & 261 \\
7.10 & 定理 & （半稳定标准模型定理）记号与假设同定理 7.9。…… & 261 \\
7.11 & 定理 & 设 $(X,B)$ 是三维 Q-分解 dlt 对，$f\colon X\to Y$ 是射影态射，而 $g\colon Y\to C$ 是到非奇异曲线的平坦态射。…… & 262 \\
7.12 & 引理 & 此时，存在 Zariski 开集 $C^\circ\subset C$，使得 $Y$ 上任一 $(K_X+B)$-翻转收缩 $h\colon X\to X^*$ 在 $(g\circ f)^{-1}(C^\circ)$ 上都是同构。 & 262 \\
7.13 & 编号段 & （定理 7.11 的证明）\quad 标准翻转的终止（定理 6.17）的证明稍加修改，便可作为定理 7.11 的证明。…… & 264 \\
7.14 & 引理 & 设 $X$ 是正规 $n$ 维簇，$S$ 是 Cartier 除子，且 $\Delta=\sum a_iD_i$，其中 $D_i$ 是互不相同的素除子，假设 $a_i\ge0$。…… & 264 \\
7.15 & 定义 & （参看定义 6.20）设 $(X,S+\Delta)$ 是 lc 对，$(X,\Delta)$ 是 dlt 对。…… & 264 \\
7.16 & 引理 & （参看引理 6.21）设 $\phi\colon X\dashrightarrow X'$ 是定理 7.11 中所述的 $(K_X+S+\Delta)$-翻转，其中 $\Delta=\sum a_iD_i$。…… & 265 \\
7.17 & 定理 & （半稳定约化定理 [KKMSD73]）记号如上。…… & 265 \\
7.18 & 注 & 若使用所需形式的终端化（定理 6.23）与 Q-分解化（定理 6.25），则下述定理 7.19 在代数设定与解析…… & 266 \\
7.19 & 定理 & （三维簇的半稳定约化定理）记号同定理 7.17。…… & 266 \\
7.20 & 注 & 定理 7.19 及其证明是二维版本 [DM69] 的自然推广。…… & 266 \\
7.21 & 编号段 & （证明的步骤与记号）\quad 设 $f'_{tf}\colon X'_{tf}\to C'$ 是诱导态射，$B'_n,B'_{tf}$ 分别是 $B$ 向 $X'_n,X'_{tf}$ 的拉回。…… & 267 \\
7.22 & 引理 & 要执行 7.21 的各阶段，只需在命题 7.23 中令 $m=3$ 所得的代数芽 $f(x_1,x_2,x_3)=\prod x_i^{a_i}$、$p(t)=t^d$ 的情形进行计算即可。 & 267 \\
7.23 & 命题 & 设 $X=(0\in\C^m)$ 是芽，$B$ 是满足 $\{\prod x_i^{a_i}=0\}\subset B\subset\{\prod x_i=0\}$ 的除子。…… & 268 \\
7.24 & 引理 & 设 $Y$ 是标准的，$D\subset Y$ 是 Cartier 除子，且假设 $(Y,D)$ 是 lc 的。…… & 269 \\
7.25 & 引理 & 按如下方式定义 $(\mu_\infty)^m$ 对命题 7.23 中引入的 $X'$ 的作用： \[ (\lambda\_1,\cdots,\lambda\_m)\colon(x\_1,\cdots,x\_m,t)\mapsto(\lambda\_1^dx\_1,\cdots,\lambda\_m^dx\_m,t\prod\lambda\_i^{a\_i}). \] 这一作用可延拓为 $X'_n$ 上的、具有有限核…… & 269 \\
7.26 & 推论 & 由引理 7.25 定义的 $X'_n$ 上的 $(\mu_\infty)^3$-作用，可延拓为 $X'_{tf}$ 上的、具有有限核的 $(\mu_\infty)^3$-作用，使 $B'^{(\ell)}_{tf}$ 是 $(\mu_\infty)^3$-不变的。 & 269 \\
7.27 & 引理 & 对任意不可约分支 $E\subset B'_{tf}$，$T_E$ 是由下式给出的 $(\mu_\infty)^3$ 的子群： \[ T\_E=\{(t^{v\_E(x\_1)},t^{v\_E(x\_2)},t^{v\_E(x\_3)})\mid t\in\mu\_\infty\}\subset(\mu\_\infty)^3. \] 其中 $v_E(\ )$ 表示沿 $E$ 的赋值。…… & 270 \\
7.28 & 引理 & 设 $(0\in Y)$ 是三维超曲面的芽，且 $Y$ 上有具有有限核的 $\mu_\infty$-作用。…… & 270 \\
7.29 & 引理 & 设 $(R,M)$ 是带有有限交换群 $G$ 作用的局部环。…… & 271 \\
7.30 & 推论 & $B'^{(\ell)}_{tf}$ 与 $X'_{tf}$ 非奇异。 & 272 \\
7.31 & 引理 & $B'_{tf}$ 在局部解析意义下是 nc 除子。 & 272 \\
7.32 & 定理 & （特殊半稳定翻转的存在）设 $(X,S+B_X)$ 是 Q-分解三维对，$f\colon X\to Y$ 是射线型 $(K+B_X)$-翻转收缩，$S$ 是 $f$-平凡的有效 Cartier 除…… & 272 \\
7.33 & 注 & 上述定理在射影的、非固有代数的以及解析的设定下都同样成立。…… & 273 \\
7.34 & 命题 & 设 $(X,S_1+S_2+B)$ 是 Q-分解三维对，$S_1,S_2$ 不可约，且 $f\colon X\to Y$ 是射线型 $(K+S_1+S_2+B)$-翻转收缩。…… & 273 \\
7.35 & 注 & $D^+$ 与 [Sho92] 中引入的 1-补足（1-complement）密切相关。…… & 273 \\
7.36 & 主张 & 设 $B^*\in|2D'|$ 是一般元。…… & 274 \\
7.37 & 注 & 在定理 7.32 中，若 $Y$ 是射影三维簇，则本节的论证只需稍作修改（不必移到开部分集合…… & 274 \\
7.38 & 命题 & 设 $T$ 是正规曲面，$p\colon S\to T$ 是奇点消解，$C$ 为例外集。…… & 276 \\
7.39 & 命题 & 设 $q\colon Y\to Z$ 是奇点消解。…… & 277 \\
7.40 & 定理 & 设 $(X,B)$ 是三维 klt 对，$B$ 是有效的。…… & 277 \\
7.41 & 注 & 本节的论证在射影的、非固有代数的以及解析的设定下都同样奏效。…… & 277 \\
7.42 & 推论 & （半稳定翻转的存在）设 $(X,B)$ 是三维 Q-分解 dlt 对，$B$ 是有效的。…… & 278 \\
7.43 & 编号段 & 定理 7.40 的证明由定理 7.44 与定理 7.51 两个主要阶段，以及把它们归约下去的容易阶段 7.54 组…… & 278 \\
7.44 & 定理 & 设 $(X,B+H)$ 是三维 Q-分解对，$B,H$ 是有效 Q-除子，且 $f\colon X\to Y$ 是射影双有理态射。…… & 278 \\
7.45 & 引理 & （定理 3.5、补足 3.6）存在满足以下 2 个条件的有理数 $\lambda\in[0,1]$。…… & 279 \\
7.46 & 引理 & $K_X+B$ 在 $Y\setminus T$ 上是 $f$-nef 的，且 $\phi$ 定义了同构态射 $X\setminus S\simeq V\setminus\phi(S)$。 & 279 \\
7.47 & 引理 & 若 $\phi$ 收缩除子 $E$，则上面的条件 (1)–(4) 在把 $f\colon X\to Y$ 换成 $V\to Y$、把 $B,S,H$ 换成 $\phi_*B,\phi_*S,\lambda\phi_*H$ 之后仍然成立。 & 279 \\
7.48 & 引理 & 若 $\phi$ 是翻转收缩，则 $\phi$ 是特殊的。…… & 279 \\
7.49 & 引理 & 若 $H\ne0$，则无论引理 7.47 与引理 7.48 的哪种情形，都能把上述手续应用于新的设定。…… & 279 \\
7.50 & 编号段 & 第二个主要阶段的主张是：若存在「好的」奇点消解，则在追加修正边界的 MMP 中，可以…… & 280 \\
7.51 & 定理 & 设 $(X,B)$ 是 $B$ 有效的三维 klt 对，$f\colon X\to Y$ 是到正规 $Y$ 的射影双有理态射，且 $s\colon Y\to C$ 是使 $s\circ f$ 为 lc 的平坦态射。…… & 280 \\
7.52 & 注 & 证明中会加入许多约化的边界成分。…… & 280 \\
7.53 & 主张 & MMP 期间，下述成立。 & 281 \\
7.54 & 编号段 & （定理 7.40 的证明）\quad 设 $a>0$ 是使 $Q=a(K_X+B)$ 成为 Cartier 的整数。…… & 283 \\
7.55 & 注 & 若 $X,Y,C$ 全是射影簇，且只证明关于射影簇的论断，则 7.54 中只有把簇收缩到邻域的讨论…… & 284 \\
7.56 & 编号段 & （模空间的紧化）\quad 长期以来已知：由所有非奇异射影簇的给定数值不变量构成的…… & 284 \\
7.57 & 定义 & 域 $k$ 上亏格 $g$ 的\textbf{稳定曲线}（stable curve）是指满足下述条件的、域 $k$ 上的 1 维固有概形 $D$：奇点至多是通…… & 285 \\
7.58 & 引理 & 设从曲面到非奇异曲线的态射 $g\colon S\to C$ 的一般纤维非奇异。…… & 285 \\
7.59 & 引理 & 设 $g\colon X\to C$ 是从簇到非奇异曲线的态射，一般纤维是标准的。…… & 286 \\
7.60 & 编号段 & 引理 7.58(1) 在高维中的类似问题仍然存在。…… & 286 \\
7.61 & 例 & 从簇到超平面截面上的锥的退化，按如下方式构成。…… & 286 \\
7.62 & 定理 & 设 $C$ 是 $\C$ 上的非奇异曲线，$C^0\subset C$ 是开子概形。…… & 288 \\
7.63 & 编号段 & （商奇点的形变空间）\quad 设 $(0\in H)$ 是奇点。…… & 288 \\
7.64 & 引理 & 设 $0\in X$ 是 3 维终端奇点，$0\in S\subset X$ 是 Cartier 除子。…… & 288 \\
7.65 & 定理 & （[KSB88, 3.9]）设 $(0\in H)$ 是曲面奇点。…… & 289 \\
7.66 & 引理 & 设 $Y$ 是正规簇，$H\subset Y$ 是 Cartier 除子，$f\colon X\to Y$ 是固有双有理态射。…… & 290 \\
7.67 & 编号段 & （曲面族的同时奇点消解）\quad 在 §4.3 中已述：Du Val 奇点的任意平坦族，至少作有限满射的基变…… & 290 \\
7.68 & 定理 & （[KSB88]）设 $S$ 是 $\C$ 上的概形，$g\colon Y\to S$ 是固有平坦态射，其任意纤维都是只有孤立奇点的曲面。…… & 290 \\
7.69 & 引理 & 设 $S$ 是固有曲面，$\bar S\to S$ 是极小奇点消解。…… & 291 \\
7.70 & 引理 & 设 $g\colon U\to V$ 是正规固有曲面之间的双有理态射。…… & 292 \\
7.71 & 编号段 & （翻转）\quad 主要结果之一，正如已多次提到的，是 3 维翻转的存在定理。…… & 292 \\
7.72 & 编号段 & （翻转与几何不变量论（GIT））\quad 自翻转被发现以来，人们已理解翻转在许多状况下出现…… & 293 \\
7.73 & 编号段 & （小平维数）\quad 设 $X$ 是非奇异射影簇。…… & 293 \\
7.74 & 编号段 & （丰度（abundance）与标准环）\quad 已在 §3.2 与 §6.1 中触及标准环的有限生成性与丰度。…… & 293 \\
7.75 & 编号段 & （藤田猜想（Fujita's conjecture））\quad 设 $X$ 是非奇异射影簇，$L$ 是 $X$ 上的丰富除子。…… & 294 \\
7.76 & 编号段 & （Fano 纤维空间的双有理几何）\quad 在 §3.8 中说明了标准模型唯一、极小模型在余维 1 意义下唯…… & 295 \\
7.77 & 编号段 & （3 维 Fano 簇）\quad MMP 促使人们详细研究其过程中产生的簇。…… & 295 \\
7.78 & 编号段 & （3 维卡拉比--丘簇）\quad 卡拉比--丘簇是标准类数值 & 295 \\
7.79 & 编号段 & （3 维椭圆簇）\quad 当簇 $X$ 具有使一般纤维为椭圆曲线的态射 $f\colon X\to Z$ 时，称 $X$ 是\textbf{椭圆的}（elliptic）。…… & 296 \\
7.80 & 编号段 & （复簇的 MMP）\quad 对固有簇或紧复簇定义 MMP 存在若干困难。…… & 296 \\
7.81 & 编号段 & （正标数或混合标数下的 MMP）\quad 人们相信 MMP 在正标数或混合标数下也起作用。 & 296 \\
7.82 & 编号段 & （多重亏格的形变不变性）\quad 最近，[Siu98] 证明了：对一般型簇的光滑族 $f\colon X\to S$，纤维的多重亏格 $h^0(X_s,\cO_{X_s}(mK_{X_s}))$ 对…… & 297 \\
\end{longtable}
\endgroup
